\documentclass{article}

\usepackage{amsmath, amssymb}
\usepackage{xcolor}
\usepackage[parfill]{parskip}
\newcommand{\red}[1]{{\color{red}#1}}
\newcommand{\dex}[1]{\frac{\partial#1}{\partial x}}
\newcommand{\dt}[1]{\frac{\partial#1}{\partial t}}
\newcommand{\ddex}[1]{\frac{\partial^2#1}{\partial x^2}}


\begin{document}

	\section{Problem statement}
	Here we try to derive and solve the master equation for the very simple
	reaction
	\begin{equation}\label{eq:reaction}
		A \longleftrightarrow B
	\end{equation}
	described by the equation
	\begin{equation}
		\left\{\begin{aligned}
			&\dot{X} = \alpha \, b - \beta a\\
			&a(t) + b(t) = N\\
			&a(0) = a_0
		\end{aligned}\right.
	\end{equation}

	with $a,b, a_0 \in \mathbb{N}$ and $a,b, a_0 < N$. This is equivalent to 
	\begin{equation}
		\left\{\begin{aligned}
			&\dot{a} = \alpha \, (N-a) - \beta \,a\\
			&a(0) = a_0
		\end{aligned}\right.
	\end{equation}

	with $a, a_0 \in \mathbb{N}, a, a_0 < N$, and we will concentrate on this latter formulation from
	here onward.

	\section{Chemical master equation}
	Following the usual methodology, we look for the chemical master reaction in
	the form
	\begin{equation*}
		\frac{dp}{dt}(t,a) = \sum_{a'} \Big[W_{aa'} \, p(t, a') - W_{a'a} \, p(t,
		a)\Big]
	\end{equation*}
	where $p(t,a)$ is the probability density function of $a$ at time $t$, and
	$W_{kl}$ is the transition rate from the state $a = l$ to the state $a=k$.
	Using the correct modeling approach \red{(to be written clearly)} one obtains that
	the probability of one transition happening during an interval $dt$ follows
	$$
		\mathbb{P}(\text{A becoming B}) = W_{ba} \, dt + o(dt)
	$$
	This allows to neglect the possibility of multiple transitions occurring in a
	single time step as their probability is of the order of $o(dt^2)$.
	Building upon this, we can write the master equation for the reaction
	\eqref{eq:reaction} as 
	\begin{equation}\label{eq:mastereq}
		\begin{aligned}
		\frac{dp}{dt}(t,a) \quad = \quad & + p(t, a-1) \, \alpha \, (N - a + 1)\\
		& + p(t, a+1) \, \beta \, (a+1)\\[6pt]
		& - p(t,a) \, \alpha \, (N-a) \\[6pt]
		& - p(t,a) \, \beta \, a
		\end{aligned}
	\end{equation}
	the first two terms describing the influx of probability from the states
	$a+1$ and $a-1$, and the last two describing the corresponding outflux.

	\section{Fokker-Planck equation}
	We now want to try and solve the equation for $p(t,a)$ via the Fokker-Planck
	method. In order to do that, we have to expand $p$ in its second argument up
	to the second order. Since $N \gg 1$ (we are talking at least hundreds of
	units), we switch our focus from $p(t,a)$ to the probability $p(t,x)$ of the
	fraction $x = \frac{a}{N}, x \in [0,1] \subset \mathbb{R}$\footnote{this is
	actually $\mathbb{Q}$, I don't know if it can be useful for something later
	maybe}\red{(also Ziv does this so it should be all pretty classic - find
	better ref tho)}. The new probability densities of interest become then: 
	\begin{align*}
		& p(t,a) \rightarrow p(t,x)\\
		& p(t,a+1) \rightarrow p(t, x + \frac{1}{N})\\
		& p(t,a-1) \rightarrow p(t, x - \frac{1}{N})
	\end{align*}
	This densities can be expanded in the second argument as:
	\begin{align*}
		& p(t,x + \frac{1}{N}) = p(t,x) + \frac{1}{N}\dex{p}(t,x) +
		\frac{1}{2N^2}\ddex{p}(t,x) + o\left(\frac{1}{N^2}\right)\\[6pt]
		& p(t,x - \frac{1}{N}) = p(t,x) - \frac{1}{N}\dex{p}(t,x) +
		\frac{1}{2N^2}\ddex{p}(t,x) + o\left(\frac{1}{N^2}\right)
	\end{align*}
	(notice the minus sign). Inserting this into \eqref{eq:mastereq} and ignoring the higher order terms,
	we can write
	\begin{equation*}
		\begin{aligned}
			\dt{p}(t,x) = &+ \left[ p(t,x) + \frac{1}{N}\dex{p}(t,x) +
			\frac{1}{2N^2}\ddex{p}(t,x) \right] \, \alpha \, (1-x+\frac{1}{N})\\
			&+ \left[ p(t,x) - \frac{1}{N}\dex{p}(t,x) +
			\frac{1}{2N^2}\ddex{p}(t,x) \right] \, \beta \, (x+\frac{1}{N})\\
			&- p(t,x) \Big[ \alpha ( 1 - x )  + \beta x \Big]
		\end{aligned}
	\end{equation*}
	and after eliminating the opposite terms, we obtain the final Fokker-Planck
	equation
	\begin{equation}\label{eq:fokker}
		\begin{aligned}
			\dt{p}(t,x) \quad = \quad & \frac{\alpha \left(1-x+\frac{1}{N}\right) +
			\beta\left(x+\frac{1}{N}\right)}{2N^2} \; \ddex{p}(t,x) \quad + \\[6pt]
			& \frac{\alpha \left(1-x+\frac{1}{N}\right) -
			\beta\left(x+\frac{1}{N}\right)}{N} \; \dex{p}(t,x) \quad + \\[6pt]
			&\frac{\alpha + \beta}{N} \; p(t,x)
		\end{aligned}
	\end{equation}
	
	\section{Solving the equation}
	Let us briefly recap the properties of the problem explored so far:
	\begin{itemize}
		\item $t \in \mathbb{R}$ is the time variable of the reaction; \item $x
		\in I = [0,1] \subset \mathbb{R}$ is the fraction of N molecules located
		in A; 
		\item $p$ is a probability density function $p : \mathbb{R}\times
		I \rightarrow [0,1]$ describing the likelihood\footnote{likelihood is
		not really correct but you get my point} of finding a fraction
		$x+dx$ of molecules in A at time $t$.
	\end{itemize}
	For ease of notation, lets define the functions
	\begin{align*}
		&P(x) = \frac{\alpha \left(1-x+\frac{1}{N}\right) +
			\beta\left(x+\frac{1}{N}\right)}{2N^2}\\[6pt]
		&Q(x) = \frac{\alpha \left(1-x+\frac{1}{N}\right) -
			\beta\left(x+\frac{1}{N}\right)}{N} \\[6pt]
		&R = \frac{\alpha + \beta}{N}
	\end{align*}
	We can then proceed to solve the FK-equation via separation of variables
	\red{(one should be able to do this as $I$ is limited and also wikipedia
	uses this scenario as an example and proceeds like this. I still don't know,
	though,	what is the creterion for picking SOV over other solution
	techniques)}. By imposing $p(t,x) = T(t) \, X(x)$, we can rewrite
	\eqref{eq:fokker} as
	\begin{equation}
		\frac{\dot{T}}{T} = P(x) \frac{X''}{X} + Q(x)\frac{X'}{X} + R = -\lambda
	\end{equation}
	with $\lambda \in \mathbb{R}^+$, and proceed to solve independently for $t$ and $x$.
	
	\subsection{Spatial part: Sturm-Liouville problem}
	The spatial part of the problem represents a \textit{classic}
	(Sturm-Liouville problem).
	We start by investigating when $P(x)$ is equal to $0$:
	\begin{equation*}
		P(\bar{x}) = \alpha\left(1-\bar{x}+\frac{1}{N}\right) + \beta
		\left(\bar{x} + \frac{1}{N}\right)
		= 0
	\end{equation*}
	By factoring out $x$ and equating, one obtains the solution as
	\begin{equation*}
		\bar{x} = \frac{(N+1) \alpha + \beta}{N (\alpha - \beta)}
	\end{equation*}
	This value of $x$ lies outside of $I$, as indeed one can verify that
	$\bar{x} > 1$ if and only if $\alpha > - (N+1) \beta$, which is trivially
	true as $\alpha, \beta, N > 0$ by definition.
	The original problem
	\begin{equation}\label{eq:clearform}
		P(x) X '' + Q(x) X' + R X = -\lambda X
	\end{equation}
	can, therefore, be put in the canonical form
	\begin{equation}\label{eq:slform}
		\frac{d}{dx}\left[p(x) \frac{dX}{dx} \right] + q(x) X = -\lambda \, w(x) X
	\end{equation}
	setting
	\begin{align*}
		p(x) = e^{\int \frac{Q}{P}}, \qquad
		q(x) = \frac{R}{P}e^{\int \frac{Q}{P}}, \qquad
		w(x) = \frac{1}{P}e^{\int \frac{Q}{P}}
	\end{align*}
	From these definitions, it can be easily seen that the coefficients indeed
	respect the necessary positivity and differentiability\footnote{actually
	this I have to check for the second derivative of P} conditions required in a
	SL-problem.
	\red{SL-theory provides an extensive corpus of methods to characterize and
	solve equations in the form \eqref{eq:slform}. I have not read this
	corpus so for now I will try with the easiest way, series expansion.}

	\paragraph{Cheating?} Before diving into the series expansion, observe that
	\eqref{eq:clearform} is naturally expressed in powers of $1/N$, the second
	derivative term representing a superior power. We could \red{(could we?)} start
	by considering only the terms of order $1/N$. The equation then becomes:
	\begin{equation}
		X'(x) + r(x) X = 0
	\end{equation}
	with
	\begin{equation*}
		r(x) = \frac{\alpha + \beta + \lambda N}{\frac{\alpha(N+1)-\beta}{N} -
		x(\alpha + \beta)}
	\end{equation*}
	Unfortunately, $r(x)$ is singular for
	\begin{equation*}
		\bar{x} = \frac{\alpha(N+1) + \beta}{N(\alpha + \beta)}
	\end{equation*}
	which happens lie in the $I$ (can be easily checked in the limit $N \to
	\infty$). 
	\red{What do we do? Do we solve separately
		\begin{align*}
			&X' + r(x) X = 0 &&\text{Far from $\bar{x}$}\\
			&P(x) X'' + (R+\lambda) X = 0 && \text{Close to $\bar{x}$}
		\end{align*}
	and glue them together in the vicinity of $\bar{x}$? I remember one can do
	something similar for the Hermite equation but I don't remember exactly how.
	And also the Hermite equation is not singular.....
	}

	\paragraph{Series solution} Due to the non-negativity of $Q(x)$ on $I$, the
	Fuchs theorem grants that the solution $X(x)$ has a power series expansion
	centered at any $x_0 \in I$ converging over the whole extent of $I$.
	We can therefore proceed by inserting the expansions
	\begin{equation*}
		X(x) = \sum_{n=0}^{\infty} c_n x^n, \quad	
		X'(x) = \sum_{n=1}^{\infty} n \, c_n x^{n-1}, \quad	
		X''(x) = \sum_{n=2}^{\infty} n \, (n-1) \, c_n x^{n-2}
	\end{equation*}
	into the original equation \eqref{eq:clearform}, obtaining
	\begin{align*}
		&\left[x\left(\frac{\beta - \alpha}{2N}\right) + \frac{\alpha(N+1) +
		\beta}{2N^2}\right]\sum_{n=2}^{\infty} n \, (n-1) \, c_n x^{n-2} \quad +
		\\
		&\left[- x\left( \beta + \alpha \right) + \frac{\alpha(N+1) -
		\beta}{N}\right]\sum_{n=2}^{\infty} n \, c_n x^{n-1} \quad + 
		\\
		&(\alpha + \beta + N\lambda) \sum_{n=0}^{\infty} c_n x^n = 0
	\end{align*}
	By rearranging the summation indexes and factoring the powers, one can
	obtain\footnote{one can but its painful}
	\begin{align*}
		&\frac{\alpha(N+1) + \beta}{N^2} \, c_2 + \frac{\alpha(N+1) -\beta}{N} \,
		c_1 + (\alpha + \beta + N\lambda) \, c_0 \quad + 
		\\ 
		&\sum_{n=1}^{\infty} 
		\Big\{\frac{\alpha(N+1) + \beta}{2N^2}(n+2)(n+1) \, c_{n+2} \quad + \\
		&\qquad \frac{n+1}{N} \left[(\beta-\alpha)\left(\frac{n}{2}-1\right) +
		\alpha N \right] \, c_{n+1} \quad + \\
		&\qquad \left[(\beta + \alpha)(1-n) + N\lambda \right] \, c_n \Big\} = 0
	\end{align*}
	\red{Appearently, each subsequent coefficient has an additional $N$ term,
		quickly exploding - can this even be well approximated by some first $n$
	series terms?}
\end{document}	
